% from ega1-10-fr.tex
\oldpage[I]{199}

faisceau d'idéaux
$\mathscr{J}'=\widetilde{\mathfrak{J}'}$ tel que le support de
$\mathscr{O}_X/\mathscr{J}'$ soit encore $X'$, les topologies
$\mathfrak{J}$-préadique et $\mathfrak{J}'$-préadique sur $A$ sont les
mêmes (\hyperref[I.10.8.2-fr]{10.8.2}).

On notera que, d'après cette définition, l'application continue
$X'\to Y'$ des espaces sous-jacents à $X_{/X'}$ et $Y_{/Y'}$, qui
correspond à $\widehat f$, n'est autre que la restriction à $X'$ de $f$.
\end{env}

\begin{env}[10.9.2]
\label{I.10.9.2-fr}
Il résulte aussitôt de la définition précédente que le diagramme de
morphismes d'espaces annelés
\[
  \label{I.10.9.2.1-fr}
  \xymatrix{
    \widehat X\ar[r]^{\widehat f}\ar[d]_{i_X} &
    \widehat Y\ar[d]^{i_Y}\\
    X\ar[r]_f &
    Y
  }
\]
est commutatif, les flèches verticales étant les morphismes canoniques
(\hyperref[I.10.8.7-fr]{10.8.7}).
\end{env}

\begin{env}[10.9.3]
\label{I.10.9.3-fr}
Soient $Z$ un troisième préschéma, $g:Y\to Z$ un morphisme, $Z'$ une
partie fermée de $Z$ telle que $g(Y')\subset Z'$. Si $\widehat g$ désigne
le complété le long de $Y'$ et $Z'$ du morphisme $g$, il résulte aussitôt
de (\hyperref[I.10.9.1-fr]{10.9.1}) que l'on a
$\widehat{(g\circ f)}=\widehat g\circ\widehat f$.
\end{env}

\begin{proposition}[10.9.4]
\label{I.10.9.4-fr}
Soient $X$, $Y$ deux $S$-préschémas localement noethériens, $Y$ étant de
type fini sur $S$. Soient $f$, $g$, deux $S$-morphismes de $X$ dans $Y$
tels que $f(X')\subset Y'$, $g(X')\subset Y'$. Pour que
$\widehat f=\widehat g$, il faut et il suffit que $f$ et $g$ coïncident
dans un voisinage de $X'$.
\end{proposition}

La condition est évidemment suffisante (sans hypothèse de finitude sur
$Y$). Pour voir qu'elle est nécessaire, remarquons d'abord que l'hypothèse
$\widehat f=\widehat g$ implique $f(x)=g(x)$ pour tout $x\in X'$. D'autre
part, la question étant locale, on peut supposer que $X$ et $Y$ sont des
voisinages ouverts affines respectifs de $x$ et de $y=f(x)=g(x)$,
d'anneaux noethériens, que $S$ est affine et que
$\Gamma(Y,\mathscr{O}_Y)$ est une $\Gamma(S,\mathscr{O}_S)$-algèbre de
type fini (\hyperref[I.6.3.3-fr]{6.3.3}). Alors $f$ et $g$ correspondent à
deux $\Gamma(S,\mathscr{O}_S)$-homomorphismes $\rho$, $\sigma$ de
$\Gamma(Y,\mathscr{O}_Y)$ dans $\Gamma(X,\mathscr{O}_X)$
(\hyperref[I.1.7.3-fr]{1.7.3}), et par hypothèse, les prolongements par
continuité de ces homomorphismes au séparé complété de
$\Gamma(Y,\mathscr{O}_Y)$ sont les mêmes. On conclut de
(\hyperref[I.10.8.11-fr]{10.8.11}) que pour toute section
$s\in\Gamma(Y,\mathscr{O}_Y)$, les sections $\rho(s)$ et $\sigma(s)$
coïncident dans un voisinage de $X'$ (dépendant de $s$)~; comme
$\Gamma(Y,\mathscr{O}_Y)$ est une algèbre de type fini sur
$\Gamma(S,\mathscr{O}_S)$, on en déduit aussitôt qu'il existe un
voisinage $V$ de $X'$ tel que $\rho(s)$ et $\sigma(s)$ coïncident dans
$V$ pour \emph{toute} section $s\in\Gamma(Y,\mathscr{O}_Y)$. Si
$h\in\Gamma(Y,\mathscr{O}_X)$ est tel que $D(h)$ soit un voisinage de $x$
contenu dans $V$, on conclut de ce qui précède et de
(\hyperref[I.1.4.1-fr]{1.4.1, d)}) que $f$ et $g$ coïncident dans $D(h)$.

\begin{proposition}[10.9.5]
\label{I.10.9.5-fr}
Sous les hypothèses de (\hyperref[I.10.9.1-fr]{10.9.1}), pour tout
$\mathscr{O}_Y$-Module cohérent $\mathscr{G}$, il existe un isomorphisme
canonique fonctoriel de $(\mathscr{O}_X)_{/X'}$-Modules
\[
  (f^*(\mathscr{G}))_{/X'}\xrightarrow{\sim}
  \widehat f^*(\mathscr{G}_{/Y'})
\]
\end{proposition}

Si on identifie canoniquement $(f^*(\mathscr{G}))_{/X'}$ à
$i_X^*(f^*(\mathscr{G}))$ et $\widehat f^*(\mathscr{G}_{/Y'})$ à
$\widehat f^*(i_Y^*(\mathscr{G}))$
(\hyperref[I.10.8.8-fr]{10.8.8}), la proposition résulte aussitôt de la
commutativité du diagramme de
(\hyperref[I.10.9.2-fr]{10.9.2}).

\begin{env}[10.9.6]
\label{I.10.9.6-fr}
Soient maintenant $\mathscr{F}$ un $\mathscr{O}_X$-Module cohérent,
$\mathscr{G}$ un $\mathscr{O}_Y$-Module cohérent. Si
$u:\mathscr{G}\to\mathscr{F}$ est un $f$-morphisme de $\mathscr{G}$ dans
$\mathscr{F}$, il lui correspond un $\mathscr{O}_X$-homomorphisme
$u^\sharp:f^*(\mathscr{G})\to\mathscr{F}$, donc par complétion un
$(\mathscr{O}_X)_{/X'}$-homomorphisme continu
$(u^\sharp)_{/X'}:(f^*(\mathscr{G}))_{/X'}\to\mathscr{F}_{/X'}$, et en
vertu de (\hyperref[I.10.9.5-fr]{10.9.5}) il existe un
$\widehat f$-morphisme $v:\mathscr{G}_{/Y'}\to\mathscr{F}_{/X'}$,
